% !TEX root = task-2-5-2.tex
% Editable task bank. Append new BankTask blocks at the end of the document.
\documentclass[12pt,a4paper]{article}
\usepackage{physics-textbook}
\usepackage{graphicx,needspace}
\usetikzlibrary{decorations.pathmorphing}
\newif\ifShowSolutions
\ifdefined\TutorEdition
  \ShowSolutionstrue
\else
  \ShowSolutionsfalse % Change to \ShowSolutionstrue to display solutions.
\fi
\newcounter{banktask}
\newsavebox{\BankTaskBox}
\newcommand{\BankTaskContent}[3]{%
  \noindent\textbf{\thebanktask.}\quad #2\par
  \ifShowSolutions
    \Needspace{50mm}% Keep the solution heading with its opening text or diagram.
    #3\par
    \ifdefined\TutorEdition
      \ifcsname TutorSource#1\endcsname
        \nopagebreak\smallskip{\small\csname TutorSource#1\endcsname\par}%
      \fi
    \fi
  \fi
}
\newcommand{\BankTask}[3]{%
  \refstepcounter{banktask}\par\medskip
  \sbox{\BankTaskBox}{\begin{minipage}{\linewidth}%
    \BankTaskContent{#1}{#2}{#3}\end{minipage}}%
  \ifdim\dimexpr\ht\BankTaskBox+\dp\BankTaskBox\relax>\dimexpr\textheight-20mm\relax
    \Needspace{35mm}\BankTaskContent{#1}{#2}{#3}%
  \else
    \Needspace{\dimexpr\ht\BankTaskBox+\dp\BankTaskBox\relax}%
    \noindent\usebox{\BankTaskBox}\par
  \fi
  \medskip
}
\begin{document}
\BookSubsection{2.5.2}{Основное уравнение МКТ идеального газа}


% Source sdamgia-803: https://phys-ege.sdamgia.ru/problem?id=803
\BankTask{sdamgia-803}{
Среднюю кинетическую энергию теплового движения молекул разреженного газа уменьшили в $2$ раза и концентрацию молекул газа уменьшили в $2$ раза. Чему равно отношение конечного давления к начальному?
}{
\smallskip\noindent\textbf{Решение}\par
Давление разреженного газа пропорционально произведению концентрации молекул газа и средней кинетической энергии теплового движения, например, для одноатомного газа: 

$p= \frac{2}{3} n\overline{E}.$

При одновременном уменьшении средней кинетической энергии теплового движения в два раза и концентрации молекул в два раза давление разреженного газа уменьшится в $4$ раза.

Ответ: $0{,}25$.
}

% Source sdamgia-811: https://phys-ege.sdamgia.ru/problem?id=811
\BankTask{sdamgia-811}{
При неизменной концентрации молекул идеального газа средняя квадратичная скорость теплового движения его молекул увеличилась в $4$ раза. Во сколько раз изменилось давление газа?
}{
\smallskip\noindent\textbf{Решение}\par
Давление идеального газа пропорционально произведению концентрации молекул газа и квадрата средней квадратичной скорости их теплового движения: 

$p= \frac{1}{3} nm_0\overline v ^{2} .$

При неизменной концентрации и увеличении средней квадратичной скорости теплового движения в $4$ раза давление газа увеличится в $16$ раз.

Ответ: $16$.
}

% Source sdamgia-3506: https://phys-ege.sdamgia.ru/problem?id=3506
\BankTask{sdamgia-3506}{
Чему равно соотношение давлений в сосудах с кислородом и водородом $p_{\text{к}}/p_{\text{в}},$ если концентрации газов и среднеквадратичные скорости одинаковы?
}{
\smallskip\noindent\textbf{Решение}\par
Давление газа пропорционально произведению его концентрации и средней кинетической энергии теплового движения: $p= \frac{2}{3} n\overline{E}.$ Последняя связана со среднеквадратичной скоростью соотношением: $\overline{E}= \frac{m v _{\text{ср.кв.}} ^{2}}{2} .$ Поскольку молекула кислорода в $16$ раз тяжелее молекулы водорода, а по условию среднеквадратичные скорости газов совпадают, средние энергии теплового движения у газов отличаются в $16$ раз. Таким образом, $p_{\text{к}}/p_{\text{в}}=16.$

Ответ: $16$.
}

% Source sdamgia-8675: https://phys-ege.sdamgia.ru/problem?id=8675
\BankTask{sdamgia-8675}{
При неизменном давлении одноатомного идеального газа среднеквадратичная скорость движения его атомов увеличилась в $2$ раза. Чему равно отношение конечной плотности газа к начальной?
}{
\smallskip\noindent\textbf{Решение}\par
Давление идеального газа можно найти по формуле $p= \frac{2}{3} nE_{\text{кин.ср.}}.$ Концентрация равна $n= \frac{N}{V} = \frac{N\rho}{m} = \frac{\rho}{m_{\text{ат}}} ,$ где $m$ — масса всего газа, $m_{\text{ат}}$ — масса атома. Средняя кинетическая энергия атомов газа $E_{\text{кин.ср.}}= \frac{m_{\text{ат}}\langle v ^{2} \rangle}{2} .$ Подставим выражение для $E_{\text{кин.ср.}}$ в уравнение для $p$ и выразим плотность газа:

$p= \frac{2}{3} \cdot \frac{\rho}{m_{\text{ат}}} \cdot \frac{m_{\text{ат}}\langle v ^{2} \rangle}{2} \Longleftrightarrow \rho= \frac{3p}{\langle v ^{2} \rangle} .$

Среднеквадратичная скорость движения атомов равна корню из среднего квадрата скорости атомов $v _{\text{кв}}= \sqrt{\langle v ^{2} \rangle} .$ Следовательно, при увеличении среднеквадратичной скорости движения атомов в $2$ раза плотность газа уменьшается в $4$ раз.

Ответ: $0{,}25$.
}

% Source sdamgia-11265: https://phys-ege.sdamgia.ru/problem?id=11265
\BankTask{sdamgia-11265}{
В закрытом сосуде находится идеальный газ при давлении $105750\,\text{Па}$ и температуре, соответствующей среднеквадратичной скорости теплового хаотического движения молекул $494\,\text{м/с}$. Чему равна плотность этого газа? (В килограммах на кубический метр; округлите до десятых долей.)
}{
\smallskip\noindent\textbf{Решение}\par
Средняя энергия теплового движения молекул связана с абсолютной температурой газа соотношением: 

$\overline{E}= \frac{m\overline v ^{2}}{2} = \frac{3}{2} kT,$ 

где $m$ — масса одной молекулы.

По определению среднеквадратичная скорость равна: 

$\bar v = \sqrt{\overline v ^{2}} = \sqrt{\frac{3kT}{m}} ,$ 

Согласно уравнению состояния идеального газа давление связано с температурой газа:

$p=nkT.$ 

Тогда преобразуем предыдущее уравнение:

$\bar v = \sqrt{\overline v ^{2}} = \sqrt{\frac{3p}{mn}} .$ 

Учитывая то, что плотность газа — это произведение массы одной молекулы на концентрацию газа, получим:

$\rho=mn= \frac{3p}{\bar v ^{2}} \approx 1{,}3\text{кг}/\text{м} ^{3} .$ 

Второй способ решения.

Основное уравнение МКТ идеального газа $p= \frac{1}{3} m_0 n \overline v ^{2} .$ В нем произведение массы молекулы на концентрацию равно $m_0 n= \frac{m}{N} \cdot \frac{N}{V} =\rho,$ поэтому $p= \frac{1}{3} \rho \overline v ^{2} ,$ откуда получается формула для нахождения плотности $\rho= \frac{3p}{\bar v ^{2}} .$

Ответ: $1{,}3$.
}

% Source sdamgia-35990: https://phys-ege.sdamgia.ru/problem?id=35990
\BankTask{sdamgia-35990}{
В сосуде находится идеальный одноатомный газ. Средняя кинетическая энергия поступательного теплового движения одной молекулы этого газа равна $40\,\text{мэВ}$, концентрация молекул равна $2{,}4\cdot 10^{25}$ $\text{м}^{-3}$. Чему равно давление газа? (В паскалях.)
}{
\smallskip\noindent\textbf{Решение}\par
Переведем среднюю кинетическую энергию в Дж: 

$40 \text{мэВ}=40 \cdot 10 ^{- 3} \cdot 1{,}6 \cdot 10 ^{- 19} =6{,}4 \cdot 10 ^{- 21} \text{Дж}.$

Давление газа найдем по формуле:

$p= \frac{2}{3} n\overline E_{\text{к}}= \frac{2 \cdot 2{,}4 \cdot 10 ^{25} \cdot 6{,}4 \cdot 10 ^{- 21}}{3} =102 400\text{Па}.$

Ответ: $102$ $400$.
}

% Source sdamgia-31885: https://phys-ege.sdamgia.ru/problem?id=31885
\BankTask{sdamgia-31885}{
В сосуде находится идеальный одноатомный газ, давление которого равно $\text{р}$. Средняя кинетическая энергия поступательного теплового движения молекул этого газа равна $\overline{E}_{\text{к}}.$ ($k$ — постоянная Больцмана) 

\TaskMatchingColumns{Физическая величина}{А)  концентрация молекул газа $n$\par Б)  температура $T$}{Утверждение}{$1$)  $\frac{2p}{3\overline{E}_{\text{к}}}$\par $2$)  $\frac{3p}{2\overline{E}_{\text{к}}}$\par $3$)  $\frac{3\overline{E}_{\text{к}}}{k}$\par $4$)  $\frac{2\overline{E}_{\text{к}}}{3k}$}

\begin{center}\begin{tabular}{|c|c|}\hline
$A$ & Б\\\hline
\strut & \strut\\\hline\end{tabular}\end{center}
}{
\smallskip\noindent\textbf{Решение}\par
А) Из формулы связи между давлением и средней кинетической энергией теплового движения молекул идеального газа $p= \frac{2}{3} n \overline{E}_{\text{к}}$ следует, что концентрация молекул газа $n= \frac{3p}{2\overline{E}_{\text{к}}}$ ($2$).

Б) Из формулы связи средней кинетической теплового движения молекул идеального газа и абсолютной температуры $\overline{E}_{\text{к}}= \frac{3}{2} kT$ следует, что температура равна $T= \frac{2\overline{E}_{\text{к}}}{3k}$ ($4$).

Ответ: $24$.
}

% Source sdamgia-9162: https://phys-ege.sdamgia.ru/problem?id=9162
\BankTask{sdamgia-9162}{
В закрытом сосуде объемом $20$ литров находится $0{,}5\,\text{моль}$ азота. Давление газа в сосуде равно $100\,\text{кПа}$. Чему равна среднеквадратичная скорость молекул этого газа? (В метрах в секунду; округлите до целого числа.)
}{
\smallskip\noindent\textbf{Решение}\par
Запишем основное уравнение МКТ газа: 

$p= \frac{1}{3} m_0n\overline v ^{2} = \frac{1}{3} \frac{m}{V} \overline v ^{2} = \frac{1}{3} \frac{\nu M}{V} \overline v ^{2} .$

Найдем среднеквадратичную скорость молекул газа:

$\overline{v}=\sqrt{\frac{3pV}{\nu M}}=\sqrt{\frac{3\cdot100\cdot10^3\cdot20\cdot10^{-3}}{0{,}5\cdot0{,}028}}\approx655\text{м/с}.$

Ответ: $655\,\text{м/с}$.

Примечание.

Азот — двухатомная молекула. Средняя кинетическая энергия, приходящаяся на $5$ степеней свободы молекул азота, равна $E = \frac{5}{2} kT.$ Однако средняя кинетическая энергия поступательного движения молекул газа — в том числе и двухатомных, и трехатомных — находится по формуле $E = \frac{3}{2} kT.$ Оставшиеся $\frac{2}{2} kT$ в данном случае соответствуют вращательным степеням свободы молекул азота.
}

% Source sdamgia-24114: https://phys-ege.sdamgia.ru/problem?id=24114
\BankTask{sdamgia-24114}{
\begin{center}\begin{tikzpicture}[task graph]
\foreach \y in {.9,1.8,2.7,3.6,4.5}{\draw[task grid] (0,\y)--(5.1,\y);}\draw[task grid] (5.1,0)--(5.1,4.5);\draw[task axis,->] (-.2,0)--(5.6,0) node[right] {$V,\,\text{м}^3$};\draw[task axis,->] (0,-.2)--(0,5.1) node[above] {$u,\,\text{м/с}$};\node[below left] at (0,0) {$0$};
\draw[task curve,domain=.18:4.5,samples=80] plot (\x,{4.5*sqrt(\x/4.5)});\foreach \x/\y/\s in {.18/.9/0{,}014,.72/1.8/0{,}056,1.62/2.7/0{,}126,2.88/3.6/0{,}224,4.5/4.5/0{,}350}{\fill (\x,\y) circle (1.6pt);}\foreach \y/\s in {.9/300,1.8/600,2.7/900,3.6/1200,4.5/1500}{\node[left] at (0,\y) {$\s$};}\foreach \x/\y/\s in {.72/1.8/0{,}056,2.88/3.6/0{,}224,4.5/4.5/0{,}350}{\draw[task axis,->] (\x,\y)--(\x,-.4);\node[below] at (\x,-.42) {$\s$};}\foreach \x/\y/\s in {.18/.9/0{,}014,1.62/2.7/0{,}126}{\draw[task axis,->] (\x,\y)--(\x,-1);\node[below] at (\x,-1.02) {$\s$};}
\end{tikzpicture}\end{center}
Два моля азота изобарно нагревают. На рисунке изображен график зависимости величины среднеквадратичной скорости $u$ молекул газа от его объема $V$. Чему равно давление газа в сосуде? (В килопаскалях.)
}{
\smallskip\noindent\textbf{Решение}\par
Воспользуемся формулами основного уравнения МКТ идеального газа, массы одной молекулы и концентрации:

$p= \frac{1}{3} m_0n\overline v ^{2} ,$ $n= \frac{N}{V} ,$ $m_o= \frac{M}{N_A} ,$ $N=\nu N_A.$ Откуда $p= \frac{\nu M\overline v ^{2}}{3V} .$

Из графика возьмем одну точку, в которой $V=0{,}014$ $\text{м}^{3}$ и $\overline v =300\text{м}/\text{с}.$ Давление газа равно: 

$p= \frac{2 \cdot 28 \cdot 10 ^{- 3} \cdot 9 \cdot 10 ^{4}}{3 \cdot 14 \cdot 10 ^{- 3}} =120\text{кПа}.$

Ответ: $120\,\text{кПа}$.
}

% Source sdamgia-29733: https://phys-ege.sdamgia.ru/problem?id=29733
\BankTask{sdamgia-29733}{
Цилиндрический сосуд разделен легким подвижным поршнем на две части. В одной части сосуда находится аргон, в другой — гелий. Концентрация атомов аргона в $2$ раза больше, чем атомов гелия. Поршень может двигаться в сосуде без трения. Определите отношение средней кинетической энергии теплового движения атома аргона к средней кинетической энергии теплового движения атома гелия при равновесии поршня.
}{
\smallskip\noindent\textbf{Решение}\par
$1$. Поскольку трение между поршнем и стенками сосуда отсутствует, то будет он перемещаться до тех пор, пока давление аргона не станет равным давлению гелия. 

$2$. Исходя из формулы связи давления от средней кинетической энергии теплового движения молекул идеального газа $p= \frac{2}{3} n \overline{E}_k,$ зависимость между средней кинетической энергией и концентрацией молекул газа обратно пропорциональная. Учитывая, что концентрация аргона в $2$ раза больше концентрации гелия, делаем вывод, что средняя кинетическая энергия молекул гелия в $2$ раза больше, чем средняя кинетическая энергия молекул аргона. 

Ответ: средняя кинетическая энергия молекул гелия в $2$ раза больше, чем средняя кинетическая энергия молекул аргона.
}


% Source 0A3E4F: https://ege.fipi.ru/bank/questions.php?search=1&proj=BA1F39653304A5B041B656915DC36B38&qid=0A3E4F
\BankTask{0A3E4F}{
Горизонтальный цилиндрический сосуд с гладкими стенками разделён подвижным поршнем на две части. В одной части сосуда находится неон, в другой — аргон. Определите отношение средних кинетических энергий теплового движения молекул неона и аргона $\frac{E_{\text{н}}}{E_{\text{а}}}$, если поршень покоится, а отношение концентраций газов $\frac{n_{\text{н}}}{n_{\text{а}}}=\frac{1}{3}$.
}{
\smallskip\noindent\textbf{Решение от GPT}\par
При равновесии поршня давления газов равны: $p_{\text{н}}=p_{\text{а}}$. По основному уравнению МКТ $p=\frac{2}{3}nE$, поэтому $n_{\text{н}}E_{\text{н}}=n_{\text{а}}E_{\text{а}}$. Следовательно, $\frac{E_{\text{н}}}{E_{\text{а}}}=\frac{n_{\text{а}}}{n_{\text{н}}}=3$. Ответ: $3$.
}

% Source 9D4F49: https://ege.fipi.ru/bank/questions.php?search=1&proj=BA1F39653304A5B041B656915DC36B38&qid=9D4F49
\BankTask{9D4F49}{
Первоначальное давление идеального одноатомного газа в сосуде равнялось \(180\,\text{кПа}\). Объём сосуда увеличили, при этом концентрация молекул газа уменьшилась в \(3\) раза. Одновременно сосуд нагрели так, что средняя кинетическая энергия теплового движения молекул газа увеличилась в \(2\) раза. Каким в результате этого стало давление \(p\) газа в сосуде (в $\text{кПа}$)?
}{
\smallskip\noindent\textbf{Решение от GPT}\par
По основному уравнению МКТ \(p=\frac{2}{3}n\overline{E}_{\text{к}}\). Тогда \(\frac{p_2}{p_1}=\frac{n_2}{n_1}\frac{\overline{E}_{\text{к},2}}{\overline{E}_{\text{к},1}}=\frac{2}{3}\), откуда \(p_2=180\cdot\frac{2}{3}=120\,\text{кПа}\).

Ответ: \(120\,\text{кПа}\).
}

% Source 3AB2B2: https://ege.fipi.ru/bank/questions.php?search=1&proj=BA1F39653304A5B041B656915DC36B38&qid=3AB2B2
\BankTask{3AB2B2}{
В ходе эксперимента давление разреженного газа в сосуде снизилось в \(5\) раз, а средняя энергия теплового движения его молекул уменьшилась в \(2\) раза. Во сколько раз уменьшилась при этом концентрация молекул газа в сосуде?
}{
\smallskip\noindent\textbf{Решение от GPT}\par
Из \(p=\frac{2}{3}n\overline{E}_{\text{к}}\) получаем \(\frac{n_2}{n_1}=\frac{p_2}{p_1}\frac{\overline{E}_{\text{к},1}}{\overline{E}_{\text{к},2}}=\frac{1}{5}\cdot2=\frac{2}{5}\). Следовательно, \(\frac{n_1}{n_2}=\frac{5}{2}=2{,}5\).

Ответ: концентрация уменьшилась в \(2{,}5\) раза.
}

% Source 772AAE: https://ege.fipi.ru/bank/questions.php?search=1&proj=BA1F39653304A5B041B656915DC36B38&qid=772AAE
\BankTask{772AAE}{
В сосуде содержится разреженный аргон, абсолютная температура которого равна \(150\ \text{К}\). Концентрацию аргона уменьшили в \(2\) раза, при этом его давление увеличилось в \(3\) раза. Определите абсолютную температуру газа в конечном равновесном состоянии. (в $\text{К}$).
}{
\smallskip\noindent\textbf{Решение от GPT}\par
Из уравнения \(p=nkT\) следует
\[
\frac{T_2}{T_1}=\frac{p_2}{p_1}\frac{n_1}{n_2}=3\cdot2=6.
\]
Следовательно, \(T_2=6\cdot150=900\ \text{К}\).
Ответ: \(900\ \text{К}\).
}

% Source 2CEC9C: https://ege.fipi.ru/bank/questions.php?search=1&proj=BA1F39653304A5B041B656915DC36B38&qid=2CEC9C
\BankTask{2CEC9C}{
Объём сосуда, содержащего \(1\,\text{моль}\) неона, уменьшили вдвое и добавили в сосуд \(1\,\text{моль}\) гелия. Температура в сосуде поддерживается постоянной.
\begin{enumerate}
\item Концентрации неона и гелия в сосуде одинаковы.
\item Внутренняя энергия неона уменьшилась.
\item Плотность газа в сосуде не изменилась.
\item Давление в сосуде увеличилось.
\item Парциальное давление неона не изменилось.
\end{enumerate}
}{
\smallskip\noindent\textbf{Решение от GPT}\par
Пусть начальные объём и давление равны \(V_0\) и \(p_0\).

1) Верно: количества неона и гелия одинаковы, они занимают общий объём, поэтому \(n_{\mathrm{Ne}}=n_{\mathrm{He}}\).

2) Неверно: \(U_{\mathrm{Ne}}=\frac32\nu_{\mathrm{Ne}}RT\) не изменяется.

3) Неверно: масса газа увеличилась, а объём уменьшился вдвое, поэтому плотность увеличилась.

4) Верно: по уравнению идеального газа \(p\propto\nu/V\), поэтому \(p_f/p_0=2\cdot2=4\).

5) Неверно: \(p_{\mathrm{Ne}}=\nu_{\mathrm{Ne}}RT/V\), поэтому парциальное давление неона увеличилось вдвое.

Ответ: \(14\).
}

% Source 559e5c: https://ege.fipi.ru/bank/questions.php?search=1&proj=BA1F39653304A5B041B656915DC36B38&qid=559e5c
\BankTask{559e5c}{
\begin{center}
\begin{tikzpicture}[>=stealth,scale=1.1]
\draw[->] (0,0)--(4,0) node[below] {$\overline{E}_{\text{к}}$};
\draw[->] (0,0)--(0,3.35) node[left] {$p$};
\node[below left] at (0,0) {$0$};
\coordinate (one) at (1.25,.7);
\coordinate (two) at (2.65,1.48);
\coordinate (three) at (2.65,3.0);
\draw[dashed] (0,0)--(one);
\draw (one)--(two)--(three);
\draw[->] (one)--(2.1,1.17);
\draw[->] (two)--(2.65,2.5);
\fill (one) circle (1.5pt) node[below right] {$1$};
\fill (two) circle (1.5pt) node[right] {$2$};
\fill (three) circle (1.5pt) node[right] {$3$};
\end{tikzpicture}
\end{center}
Один моль аргона участвует в процессе $1$--$2$--$3$, представленном на графике зависимости давления $p$ от средней кинетической энергии $\overline{E}_{\text{к}}$ теплового движения молекул.
\begin{enumerate}
\item В процессе $1$--$2$ концентрация аргона остаётся неизменной.
\item В процессе $1$--$2$ аргон отдаёт положительное количество теплоты.
\item В процессе $2$--$3$ аргон изотермически расширяется.
\item В процессе $2$--$3$ аргон отдаёт окружающей среде положительное количество теплоты.
\item В процессе $2$--$3$ температура аргона растёт.
\end{enumerate}
}{
\smallskip\noindent\textbf{Решение от GPT}\par
Используем $p=\frac23 n\overline{E}_{\text{к}}$ и $\overline{E}_{\text{к}}=\frac32 kT$.
\begin{enumerate}
\item Верно: на $1$--$2$ отношение $p/\overline{E}_{\text{к}}$ постоянно, следовательно, концентрация $n$ постоянна.
\item Неверно: $1$--$2$ --- изохорное нагревание, $Q=\Delta U>0$, газ получает теплоту.
\item Неверно: на $2$--$3$ температура постоянна, давление растёт, поэтому объём уменьшается: происходит сжатие.
\item Верно: при изотермическом сжатии $\Delta U=0$, работа газа $A<0$, поэтому $Q=A<0$.
\item Неверно: на $2$--$3$ средняя кинетическая энергия и температура постоянны.
\end{enumerate}
\textbf{Ответ: 14.}
}

% Source 5644DE: https://ege.fipi.ru/bank/questions.php?search=1&proj=BA1F39653304A5B041B656915DC36B38&qid=5644DE
\BankTask{5644DE}{
В цилиндре под поршнем находится идеальный одноатомный газ. Формулы А и Б позволяют рассчитать значения физических величин, характеризующих состояние газа. Использованы обозначения: \(p\) — давление; \(V\) — объём; \(\overline{E}_{\text{к}}\) — средняя кинетическая энергия поступательного теплового движения частиц газа; \(\nu\) — количество вещества, \(N_A\) — постоянная Авогадро.

Установите соответствие между формулами и физическими величинами, значение которых можно рассчитать по этим формулам.

\begin{center}
\TaskMatchingColumns[.35\linewidth]{Формулы}{
А) \(\displaystyle\frac{2}{3}\frac{\nu N_A\overline{E}_{\text{к}}}{V}\)\par
Б) \(\displaystyle\frac{3}{2}\frac{pV}{N_A\overline{E}_{\text{к}}}\)
}{Физические величины}{
1) масса газа\par
2) количество вещества\par
3) давление\par
4) абсолютная температура
}

\begin{tabular}{|c|c|}
\hline
А & Б \\
\hline
\phantom{0} & \phantom{0} \\
\hline
\end{tabular}
\end{center}
}{
\smallskip\noindent\textbf{Решение от GPT}\par
Основное уравнение МКТ: \(pV=\frac{2}{3}N\overline{E}_{\text{к}}\), где \(N=\nu N_A\). Поэтому \(p=\frac{2}{3}\frac{\nu N_A\overline{E}_{\text{к}}}{V}\): А — 3. Выражая количество вещества, получаем \(\nu=\frac{3}{2}\frac{pV}{N_A\overline{E}_{\text{к}}}\): Б — 2.

Ответ: А — 3, Б — 2; \(32\).
}

% Source 0D4FC9: https://ege.fipi.ru/bank/questions.php?search=1&proj=BA1F39653304A5B041B656915DC36B38&qid=0D4FC9
\BankTask{0D4FC9}{
При исследовании изопроцессов использовался закрытый сосуд переменного объёма, заполненный неоном и соединённый с манометром. Объём сосуда медленно увеличивают, сохраняя давление неона в нём постоянным. Как изменяются при этом внутренняя энергия и плотность неона в сосуде?

Для каждой величины определите соответствующий характер её изменения:
\begin{enumerate}
\item увеличивается
\item уменьшается
\item не изменяется
\end{enumerate}

\TaskChangeTable{Внутренняя энергия неона в сосуде}{Плотность неона в сосуде}
}{
\smallskip\noindent\textbf{Решение от GPT}\par
Для одноатомного идеального газа \(U=\frac32\nu RT=\frac32pV\). При постоянном давлении увеличение объёма приводит к увеличению внутренней энергии (1). Масса газа в закрытом сосуде постоянна, поэтому плотность \(\rho=m/V\) уменьшается (2).

\(\boxed{12}\).
}

\end{document}
